\section{Discussion: Policy Design}\label{sec:discussion}

The policy lesson is not that set-asides are universally dominated. It is that bidder-exclusion rules should be evaluated against the thickness of the protected pool, the heterogeneity of the good, and the welfare weight the planner places on SME surplus. Law 14{,}133/2021 allows both full SME set-asides and price preferences. The estimates here imply that the choice between them should be contingent on market structure rather than imposed as a uniform default.

\subsection{When does a set-aside become defensible?}

The structural decomposition suggests a simple answer. A full set-aside becomes more defensible when three conditions line up: the SME pool is thin enough that protecting it is a first-order objective, the good is heterogeneous enough that the identity of the induced entrants matters for equilibrium selection, and the planner places a high welfare weight on SME surplus. Those are precisely the conditions under which the policy ranking becomes model-sensitive in the data.

Non-pharma sits at the opposite corner. The baseline SME pool is thicker, the goods are more standardized, and the 10 percent price preference is Pareto-superior under both the main equilibrium-selection reading and the strict-invariance benchmark. That is the market class in which bidder exclusion is hardest to defend. Pharma is closer to the knife-edge: the pool is thinner, heterogeneity is higher, and the policy ranking depends on whether the post-policy entrant mix is treated as selected equilibrium participation or forced to equal the pre-policy primitive. The point is not that pharma is special. The point is that thin, heterogeneous markets are exactly where bidder-exclusion rules should be expected to be more sensitive.

\subsection{The welfare weight implicit in the current regime}

Consider a social planner aggregating producer and consumer surplus with weight $w^{\text{SME}}$ on SME producer surplus. The planner is indifferent between the full set-aside (V0) and the 10 percent price preference (V3) when
\begin{equation}
w^{\text{SME}}_\star
\;=\; \frac{\mathrm{Loss}(V_0) - \mathrm{Loss}(V_3)}
           {\mathrm{SMESurplus}(V_0) - \mathrm{SMESurplus}(V_3)}.
\end{equation}
Under the main specification, the implied indifference weight is about 2.4 in non-pharma and 3.0 in pharma. Under the strict-invariance benchmark, the non-pharma preference ranking is unchanged, while the pharma indifference weight falls to 0.7. These are not contradictory numbers. They are a map of when bidder exclusion is doing something that a moderate preference rule cannot.

The comparison is informative because it disciplines the distributional argument in favor of set-asides. A planner willing to pay the non-pharma fiscal cost of bidder exclusion would have to value one real of SME surplus at more than twice a real of general welfare, even in a thick market where a preference rule achieves nearly the same allocative objective. In pharma, the required weight depends on how strongly one believes the policy selects a different entrant mix. The right reading is therefore conditional: bidder exclusion is hardest to justify in thick standardized markets, and most plausible only in thin heterogeneous markets with very high welfare weights on protected firms. Figure~\ref{fig:v3_welfare_weight} in Section~\ref{sec:pareto} plots the implied weights against the Brazilian social-policy benchmark.

\subsection{Implementation considerations}

Two features qualify the policy comparison. First, V3 is computed in a Vickrey-equivalent framework. In a first-price sealed-bid auction with the same preference rule, bidder equilibrium behavior would adjust in ways characterized by \citet{maskinriley2000}. The direction of the adjustment typically preserves the main insight here: keeping non-SMEs active in the auction is what disciplines the price distribution. What would move is the exact surplus split.

Second, a preference rule is procedurally more complex to implement than a set-aside. The buyer must score bids twice and disclose the winner-selection rule clearly. But this is an operational complication, not a conceptual obstacle. Preference systems have been run elsewhere for decades, and Brazilian procurement law already contemplates them. The relevant policy question is therefore not whether the instrument exists, but in which markets it should replace bidder exclusion.

\subsection{Market thickness and goods heterogeneity}

The non-pharma versus pharma comparison is best read as a statement about the interaction between bidder exclusion and goods characteristics. Pharma auctions have both a thinner SME pool and higher within-class heterogeneity. Once non-SMEs are removed from price formation, those two features raise the expected second-order statistic and make the selected entrant mix more important. Non-pharma auctions are closer to the textbook standardized-good environment in which a moderate preference rule can tilt allocation toward SMEs without materially changing price formation.

This generalizes beyond the Group-65 split. Goods with deeper supplier bases and tighter specification standardization should be natural candidates for preference rules rather than exclusion. Goods with thinner protected pools, more heterogeneous production technologies, and stronger administrative reasons to value SME surplus are the cases in which bidder exclusion becomes more arguable. The contribution of the paper is to put that intuition on structural objects: the size of the intensive margin, the magnitude of the entry insurance, and the welfare weight needed to justify the policy.

\subsection{Two objections worth addressing}

Two lines of skepticism about the structural approach deserve a direct answer.

\paragraph{IPV asymmetric versus contract specificity.} A reader might argue that pharma procurement involves item-level specifications---ANVISA compliance, brand-name regulatory history, storage chain, production line provenance---that are not well captured by an independent-private-values model with only type-level asymmetry. The concern is that what looks like a bidder-specific private value is in fact a common quality signal that correlates bidder valuations within an auction, in which case Assumption~\ref{a:ipv} is misspecified and the estimated $F_c^k$ absorbs variation that should be attributed to a common factor.

Three features of the BEC setting address this. First, the sample is restricted to Preg\~ao, the electronic reverse-auction format reserved for goods with standardized specifications at the item level; BEC's CADMAT catalog is designed precisely to strip cross-item quality variation within class. This is a different data-generating environment from the construction contracts studied in \citet{bajari2014} or the Bayesian auctions of \citet{hendricks2003}, where contract-specific complexity drives correlation. Second, the CMED price ceiling applied to pharma class 6531 is itself a standardization device: any CMED-regulated drug clears a uniform quality bar, and the residual variation across suppliers is on production cost and distribution scale rather than on quality primitives. Third, the Krasnokutskaya auction-level heterogeneity correction (Assumption~\ref{a:uh}) allows for an auction-common component $a_t$ that absorbs exactly the kind of unobserved scale or specification variation the objection raises. The Turnbull NPMLE robustness relaxes Assumption~\ref{a:uh} entirely and delivers the same decomposition within five percent, so the main results do not hinge on the Gaussianity of $a_t$ or on the correctness of the BLP shrinkage.

\paragraph{Policy correlation at the cutoff.} A second concern is that the March 2018 cutoff might coincide with other Brazilian regulatory events that confound identification---CMED adjustments, platform upgrades, tax changes, federal procurement reform. Each is worth checking explicitly.

CMED list-price adjustments occur annually at the beginning of each year and apply to all CMED-regulated drugs, not to Group 65 specifically; a CMED adjustment would appear in the pre-period Group-65 data as well as in the never-treated product groups used as controls, and the DiD of Section~\ref{sec:rf} would difference it out. BEC platform architecture was stable throughout the sample window; the catalog and auction format did not change. Federal Law 14{,}133/2021 did not exist during the window---the relevant federal framework was the 2006 statute, unchanged throughout. Any Group-65-specific shock that could substitute for the SME rule would need to coincide with March 2018, hit only Group 65, and leave no trace in publicly available administrative records. Section~\ref{sec:rf} reports the pre-trend and placebo diagnostics directly, and finds no anomalous Group 65 movement before the cutoff. The cutoff is as close to exogenous as a legalistic reinterpretation inside the state prosecutor's office can be.

\subsection{External validity: the portability of the decomposition framework}\label{sec:external}

The paper studies a specific Brazilian reform, but the decomposition framework applies to every procurement rule that restricts the bidder pool on an identity criterion. SME set-asides are the most obvious case and the most widely deployed; the U.S. Small Business Administration has operated them since 1953, the European Union catalogues them in more than one hundred jurisdictions \citep{bosio2022}, and Brazilian Law~14{,}133/2021 made them a procurement default. But the same two-margin structure also characterizes rules that restrict the bidder pool by geography (Buy American Act, local-content requirements under WTO Annex~IV~Article~XXIII, state-level in-state-bidder preferences), by production method (environmental or labor-standard screens, ESG-compliance thresholds in European tenders), by national security clearance (ITAR and defense-procurement restrictions), or by historical disadvantage (minority-owned business set-asides under 13~CFR~124, women-owned business preferences under 13~CFR~127).

Appendix Table~\ref{tab:v3_external_policies} organizes these rules by mechanism. Full-exclusion rules (Block~A)---the paper's SBA set-asides, the US Berry Amendment, Brazilian Lei Complementar 123/2006 and Law 14{,}133/2021, and the S\~ao Paulo rule studied here---are precisely the regimes to which the sheltered-bidding decomposition applies without additional parametric structure. Price-preference rules (Block~B)---the HUBZone program, the Buy American Act, Italian Code 50/2016---retain both bidder types in the auction and therefore entangle the intensive and extensive margins in the way \citet{marion2007}, \citet{krasnokutskaya2011}, and \citet{atheyseira2011} address with parametric bid-function inversion. Target and liberalization rules (Block~C)---EU Directive 2014/24, the WTO GPA, the US Trade Agreements Act---operate through aggregate volumes or country origins rather than per-auction eligibility, and the decomposition does not apply.

In every full-exclusion rule the policy question the decomposition answers is the same: how much of the price cost flows through the set of admissible firms versus how each remaining firm bids? The identification does require a clean natural-experiment trigger and access to bidder-level drop-out data; the first is the rare object and the second is increasingly available as national procurement platforms digitalize. Future applications of the sheltered-bidding concept to U.S. SBA small-business set-asides, the Berry Amendment's industrial-specialty carve-outs, and local-content rules in emerging-market procurement are natural extensions. Whether the 75--85 percent intensive-margin share is specifically Brazilian or general to full-exclusion rules is an open question this paper is not positioned to answer, but is positioned to motivate.

\subsection{Static welfare cost as a lower bound on long-run consequence}\label{sec:dynamics}

The model treats procurement quantity $Q$ as inelastic and the set-aside as static. Two dynamic margins, both absent from the welfare arithmetic, deserve explicit attention because their omission affects whether the static estimate is best read as a lower or upper bound on the policy's long-run cost.

The first is the buyer's quantity response. If the SME-only rule raises unit prices and the buyer's procurement budget is capacity-constrained, fewer units are purchased per real of public spending. In the medical-supplies setting studied here, a reduced supply of consumables to the SES network translates into directly worse health outcomes, which the static welfare arithmetic does not price. The static welfare cost of 29--45 percent of procurement value is therefore a \emph{lower bound} on the social cost of the rule under any specification in which procurement quantity affects welfare.

The second is the dynamic capacity response of SMEs to a protected market. \citet{szerman2023} and \citet{cabral2017} document that protected procurement allocations can build firm-level capacity that persists beyond the protection window---the policy's intended distributive payoff. If long-run capacity growth among SMEs is sufficient, the policy's static welfare cost would be partially offset by a dynamic SME productivity gain. But two features of the present setting limit the offset. First, the realized incidence (Section~\ref{sec:welfare}) shows the set-aside's transfer is concentrated in micro-employer firms that were already in the supplier registry pre-policy; the policy is rewarding incumbents rather than financing the entry of new SMEs. Second, the documented Brazilian dynamic effects \citep{cabral2017} are calibrated on programs that combine protection with technical assistance and credit access; the BEC set-aside studied here provides protection without the complementary inputs. The dynamic offset, if any, is therefore unlikely to overturn the static ranking. The static welfare cost should be read as a credible lower bound on the long-run social cost, with a margin of upward revision that grows in the size of the dynamic capacity response.

A full dynamic-general-equilibrium treatment of the SME-only rule, with explicit modeling of the quantity response and SME productivity dynamics, is a paper of its own.
